\Askhsh{1874}{4}
\begin{ylh}{1}
    \item 3.2. 1ο Κριτήριο ισότητας τριγώνων
    \item 4.6. Άθροισμα γωνιών τριγώνου
    \item 5.9. Μια ιδιότητα του ορθογώνιου τριγώνου.
\end{ylh}
Έστω κύκλος με κέντρο $\mathrm{O}$ και διάμετρο $\mathrm{K}\Lambda$. Έστω $\mathrm{A}$ σημείο του κύκλου ώστε η ακτίνα $\mathrm{O}\mathrm{A}$ να είναι κάθετη στην $\mathrm{K}\Lambda$. Φέρουμε τις χορδές $\mathrm{A}\mathrm{B} = \mathrm{A}\Gamma = \rho$. Έστω $\Delta$ και $\mathrm{E}$ τα σημεία τομής των προεκτάσεων των $\mathrm{A}\mathrm{B}$ και $\mathrm{A}\Gamma$ αντίστοιχα με την ευθεία της διαμέτρου $\mathrm{K}\Lambda$.
Να αποδείξετε ότι:
\begin{alist}
    \item Η γωνία $\mathrm{B}\hat{\mathrm{A}}\Gamma$ είναι $120^\circ$. \monades{7}
    \item Τα σημεία $\mathrm{B}$ και $\Gamma$ είναι μέσα των $\mathrm{A}\Delta$ και $\mathrm{A}\mathrm{E}$ αντίστοιχα. \monades{9}
    \item $\mathrm{K}\Gamma = \Lambda\mathrm{B}$. \monades{9}
\end{alist}
